\documentclass[conference]{IEEEtran}

\usepackage{amsmath,amssymb}
\usepackage{booktabs}
\usepackage{graphicx}
\usepackage[hidelinks]{hyperref}
\usepackage{xcolor}

\definecolor{implementedgreen}{RGB}{0,110,55}
\definecolor{plannedblue}{RGB}{30,80,160}
\definecolor{openorange}{RGB}{180,95,0}
\newcommand{\implemented}{\textcolor{implementedgreen}{\textbf{Implemented}}}
\newcommand{\planned}{\textcolor{plannedblue}{\textbf{Planned}}}
\newcommand{\openitem}{\textcolor{openorange}{\textbf{Open}}}

\title{Trajectory-Aware Control for Autonomous\\
Tractor--Semitrailer Reverse Docking}
\author{
  \IEEEauthorblockN{Nandan Prince}
  \IEEEauthorblockA{MSML 642}
  \and
  \IEEEauthorblockN{Chris Pedretti}
  \IEEEauthorblockA{MSML 642}
}

\begin{document}

\maketitle

\begin{abstract}
This report documents the staged development of a simulation-based autonomous
reverse-docking system for a tractor connected to one conventional
semitrailer. The core study will compare a trajectory-relative Proximal Policy
Optimization (PPO) controller with a nonlinear model predictive controller
(NMPC) under a shared vehicle, task, safety, and evaluation definition. A fast
two-dimensional kinematic implementation is the required experimental
platform. The shared vehicle core and a configuration-generated MuJoCo vehicle,
dock, and adapter are implemented and covered by engineering validation. The
canonical two-dimensional propagation, Gymnasium task, controllers, direct
cross-simulator parity, training, and comparative results are not yet
implemented and are explicitly identified as planned work.
\end{abstract}

\begin{IEEEkeywords}
tractor--semitrailer, reverse docking, reinforcement learning, proximal policy
optimization, nonlinear model predictive control, MuJoCo
\end{IEEEkeywords}

\section{Introduction}

Reversing an articulated tractor--semitrailer is a nonlinear and nonholonomic
control problem. In reverse motion, small steering errors can increase the
relative heading between the tractor and trailer, causing tracking failure,
collision, or jackknifing. Loading-dock operation adds tight lateral clearance
and final-pose requirements, making it a useful setting for comparing learned
and model-based control.

The project studies whether a PPO policy using trajectory-relative tracking
errors and preview information can reliably execute reverse docking and how its
behavior compares with NMPC when initial pose, trailer geometry, parking-bay
dimensions, and state or model uncertainty change. The comparison is not
designed to make one controller win. It is intended to characterize accuracy,
safety, computation, robustness, and transfer behavior under a common task
definition.

The required contribution of the course project is an implemented and
quantitatively evaluated two-dimensional system. A three-dimensional MuJoCo
model is a secondary integration target. If full transfer is not reliable or
cannot be completed within the schedule, the complete controller comparison
will remain in the verified two-dimensional environment, while MuJoCo will be
used for a smaller validated demonstration set.

This is a living report. Statements are limited to implemented behavior and
saved evidence. Planned methods and experiments are labeled as such, and no
novelty or performance claim is made before comparative results exist.

\section{Related Work}

Yan et al.\ combine a low-speed tractor--trailer kinematic model, reference
trajectory errors, preview states, and PPO for reverse path following
\cite{yan2024trajectory}. Their parameterization includes tractor wheelbase,
rear-axle-to-hitch offset, and hitch-to-trailer-axle distance. Their work
motivates the implemented trajectory-relative observation in this project, but its
software interface and several reproduction details are not adopted verbatim.

Moradi and Mehrandezh formulate an NMPC local planner for forward and reverse
tractor--trailer motion, including state and control constraints, articulation
limits, obstacle avoidance, and a transition from a headless Python model to a
Gazebo/ROS implementation \cite{moradi2022nmpc}. Their compact experimental
dimensions and staged simulator workflow motivate the nominal vehicle scale
and development sequence used here. Our primary comparison differs by using a
shared fixed-speed, steering-rate plant interface for PPO and NMPC.

Ma et al.\ separate long-horizon reference-line graph search from constrained
reinforcement-learning trajectory generation \cite{ma2025hierarchical}. This
supports a modular separation between path generation and local execution, but
multi-lane yard routing is outside the critical path of the present study.
Menz and Wenzel study PPO-generated demonstrations, behavioral cloning, and
offline reinforcement learning for a steerable drawbar-trailer configuration
\cite{menz2026imitation}. Their experiment-management and robustness questions
are relevant, although their vehicle mechanism differs from the conventional
semitrailer used here.

PPO provides a continuous-action policy-gradient baseline with clipped policy
updates \cite{schulman2017ppo}. NMPC will be implemented using CasADi
\cite{andersson2019casadi} and an interior-point nonlinear-programming solver.
Stable-Baselines3 is the planned PPO implementation \cite{raffin2021sb3}, and
MuJoCo is the implemented rigid-body physics backend \cite{todorov2012mujoco}.

\section{Problem Formulation}

\subsection{Task}

The task contains one front-steered tractor and one non-steered, single-axle
semitrailer operating at low speed on a level surface. The map and target dock
are known. Simulator ground truth supplies the core state; perception, mapping,
and dock recognition are not required for the primary comparison.

An episode succeeds when the trailer docking point reaches the target within
longitudinal, lateral, and heading tolerances, the complete trailer footprint
is inside the bay, no collision has occurred, and articulation remains below
the hard safety limit. Collision, jackknife, operating-boundary exit, numerical
failure, and timeout will be categorized separately.

\subsection{Coordinate conventions}

The world frame is right-handed. The planar $x$ axis points right in an
overhead plot, $y$ points up, and MuJoCo $z$ points upward. Body $x$ points
forward and body $y$ points left. Yaw is counterclockwise from world $+x$.
Longitudinal speed is positive forward and negative in reverse. Steering is
positive to the left.

The tractor pose $(x_t,y_t,\psi_t)$ is measured at the tractor rear-axle
midpoint. The trailer pose $(x_r,y_r,\psi_r)$ is measured at the trailer-axle
midpoint. Articulation is
\begin{equation}
  \alpha = \operatorname{wrap}(\psi_t-\psi_r),
\end{equation}
with values in $[-\pi,\pi)$.

\subsection{Shared state and command}

The implemented simulator-independent state contract is
\begin{align}
  s = (&x_t,y_t,\psi_t,x_r,y_r,\psi_r,x_h,y_h,\delta,\dot{\delta},v,\\
       &\dot{\psi}_t,\dot{\psi}_r,t),
\end{align}
where $(x_h,y_h)$ is the measured hitch position, $\delta$ is the virtual
bicycle steering angle, $v$ is longitudinal speed at the tractor rear axle,
and $t$ is simulation time. Articulation and articulation rate are derived:
\begin{equation}
  \dot{\alpha}=\dot{\psi}_t-\dot{\psi}_r.
\end{equation}
They cannot be independently assigned in the shared state object.

The high-level plant command is
\begin{equation}
  u=(\dot{\delta}_{\mathrm{cmd}},v_{\mathrm{target}}).
\end{equation}
For the primary comparison, $v_{\mathrm{target}}=-0.2\,\mathrm{m/s}$ is fixed
by configuration and only steering rate is selected by the controller. The
target-speed field is retained for scripted validation and separately defined
future experiments.

\section{System Design and Current Implementation}

\subsection{Versioned vehicle and dock}

The initial cross-simulator specification is
\texttt{compact-semitrailer-v1}. It is a compact research-scale configuration
anchored to the NMPC paper's reported experimental dimensions, rather than a
claim of full-size highway-vehicle fidelity. Table~\ref{tab:geometry} summarizes
the frozen nominal planar geometry.

\begin{table}[ht]
  \centering
  \caption{Nominal version-1 geometry.}
  \label{tab:geometry}
  \begin{tabular}{lr}
    \toprule
    Parameter & Value \\
    \midrule
    Tractor wheelbase & $1.90\,\mathrm{m}$ \\
    Tractor body length $\times$ width & $2.90\times1.00\,\mathrm{m}$ \\
    Rear-axle-to-hitch offset & $0.00\,\mathrm{m}$ \\
    Hitch-to-trailer-axle distance & $4.00\,\mathrm{m}$ \\
    Trailer body length $\times$ width & $5.50\times1.00\,\mathrm{m}$ \\
    Wheel radius & $0.30\,\mathrm{m}$ \\
    \bottomrule
  \end{tabular}
\end{table}

The hitch offset is signed and configurable even though the nominal vehicle is
on-axle. Trailer variants use hitch-to-axle distances of $2$, $4$, and
$5\,\mathrm{m}$. An optional full-size interactive configuration now
demonstrates the same schema without replacing the compact experimental
baseline. It uses a $5.8\,\mathrm{m}$ tractor and a
$14.0\times2.55\,\mathrm{m}$ semitrailer in a $200\times100\,\mathrm{m}$ 2D
service-center yard with 20 doors along the bottom building edge and one active
$18\,\mathrm{m}$ loading zone. The target zone is highlighted by floor color
and deliberately has no physical or collidable longitudinal side walls; its
width remains part of the success containment test. The rear boundary and
building façade outside the active door are collision geometry. Faster
teleoperation limits make manual maneuvering observable at this scale; these
are demonstration choices rather than identified physical parameters.

The nominal dock interior is $6.00\times1.50\,\mathrm{m}$. Its target is the
ground-plane projection of the trailer rear-face center at
$(0.25,0.00)\,\mathrm{m}$ in the dock frame, with zero trailer yaw. Position
tolerances are $0.10\,\mathrm{m}$ longitudinally and laterally; heading
tolerance is $3^\circ$.

\subsection{Limits and timing}

The nominal steering-angle limit is $0.50\,\mathrm{rad}$ and the selected
steering-rate limit is $0.30\,\mathrm{rad/s}$. The latter is a project design
choice and will require sensitivity analysis; it is not reported by either core
paper. The soft and hard articulation thresholds are $35^\circ$ and $45^\circ$.
These values describe the compact experimental baseline. The separate
full-size interactive service-center scenario retains the $35^\circ$ warning
but terminates at $60^\circ$ articulation.
The controller period is $0.10\,\mathrm{s}$, the implemented 2D integration step is
$0.02\,\mathrm{s}$, and the MuJoCo physics step is $0.002\,\mathrm{s}$.

\subsection{Service-center presentation frame}

The service-center task retains the dock-local kinematic state while presenting
the building along the bottom edge of a conventional yard plot. If
$(x_d,y_d,\psi_d)$ is a pose in the dock frame, the renderer applies
\begin{equation}
  x_s=100-y_d,\qquad y_s=x_d,\qquad
  \psi_s=\psi_d+\frac{\pi}{2}.
\end{equation}
This transformation is a rigid $90^\circ$ rotation followed by translation,
so it preserves distances and hitch coincidence. An articulated-state
regression test separately transforms the tractor pose, trailer pose, and hitch
point and requires both transformed body hitch locations to agree with the
transformed hitch marker to numerical precision.

\subsection{Implemented software contracts}

\implemented{} components currently include immutable state and command
objects, strict vehicle and dock configuration loaders, geometry and frame
transforms, trailer variants, semantic-site definitions, state serialization,
SHA-256 configuration fingerprints, geometry-residual diagnostics, command
supervision, and an executable conformance report. Unknown, missing,
non-finite, and cross-field-inconsistent configuration values are rejected.

The 2D and MuJoCo implementations will share semantic frames for tractor rear
axle, tractor front axle, hitch, trailer axle, trailer rear face, and docking
point. MuJoCo body origins may differ, but named sites and the adapter must
recover the shared state convention.

The versioned vehicle configuration now also freezes the complete MuJoCo
identifier contract for the tractor and trailer bodies, steering links, wheel
bodies, free joint, hitch and steering joints, passive wheel-spin joints, and
four actuators. This naming layer is part of the simulator API: it permits the
implemented adapter and parity tests to resolve model elements deterministically
without depending on XML ordering.

A first MuJoCo model now implements this contract. The MJCF for the vehicle
and the nominal dock is generated from the same YAML files that define the
planar geometry, and an adapter returns the shared planar state from the named
sites at the $10\,\mathrm{Hz}$ command boundary. Wheels contact the ground
through a thin cylinder on the wheel centre plane, because full-width rigid
cylinders scrubbed at both rims and saturated the steering actuators. On a
scripted $3\,\mathrm{m}$ low-speed maneuver the model ended within
$0.009\,\mathrm{m}$ and $0.001\,\mathrm{rad}$ of a numerical integration of
the on-axle kinematic equations. This is a model-validation check on one
scenario. Direct tests now compare the project's 2D and MuJoCo backends under
shared forward and reverse command sequences; task-event and controller
transfer checks remain planned.

\subsection{Reproducibility}

Vehicle and dock YAML files contain separate schema versions and specification
identifiers. Normalized configuration fingerprints are designed to be recorded
with future experiment manifests. A new dimension set therefore changes the
result provenance without requiring controller-interface changes.

MuJoCo-only gains, damping, armature, contact width, and solver choices are
versioned in code rather than in the version-1 YAML schema. Consequently, a
reported MuJoCo experiment must also record the source revision and resolved
tuning values; the vehicle fingerprint alone is insufficient provenance.

\section{Methods}

This section distinguishes implemented infrastructure from planned control
methods.

\subsection{Two-dimensional vehicle model}

\implemented{} The planar plant uses a low-speed kinematic bicycle for the
tractor and a no-lateral-slip single-axle trailer. Its independent state is
$[x_t,y_t,\psi_t,\psi_r,\delta,v]$. With wheelbase $L_t$, trailer length $L_r$,
signed longitudinal and lateral hitch offsets $a$ and $b$, tractor yaw rate
$\omega=v\tan(\delta)/L_t$, and articulation
$\alpha=\psi_t-\psi_r$, the trailer equation is
\begin{equation}
 \dot{\psi}_r = \frac{(v-b\omega)\sin\alpha+a\omega\cos\alpha}{L_r}.
\end{equation}
Together with $\dot{x}_t=v\cos\psi_t$, $\dot{y}_t=v\sin\psi_t$,
$\dot{\psi}_t=\omega$, $\dot{\delta}=u_\delta$, and bounded acceleration
$\dot v=u_v$, this reduces to the frozen on-axle model when $a=b=0$.

The simulator applies command supervision at $10\,\mathrm{Hz}$ and advances
five fixed $0.02\,\mathrm{s}$ classical fourth-order Runge--Kutta substeps.
Hitch and trailer-axle positions are reconstructed from rigid geometry rather
than independently integrated. Tests cover analytic derivatives, straight and
turning motion in both speed directions, articulation, angle wrapping, limits,
rigid reconstruction, deterministic replay, timestep convergence, and direct
low-speed trajectory comparison with MuJoCo.

\subsection{Reference trajectories and observations}

\implemented{} Reference paths are non-degenerate polylines ordered from the
dock target toward the start so a reverse maneuver moves toward decreasing arc
length while the nominal vehicle heading remains outward. Orthogonal projection
provides signed lateral error and progress; preview points are sampled toward
the goal at $0.5$, $1.0$, and $2.0\,\mathrm{m}$.

The bounded 21-element observation contains lateral error, sine/cosine tractor
and trailer heading errors, articulation and articulation rate, steering state,
speed, remaining path distance, docking-pose error, and the local coordinates
of the three preview points. Values are normalized and clipped to $[-1,1]$,
with a stable name assigned to every element.

\subsection{Gymnasium docking task}

\implemented{} The registered environment advances one $0.10\,\mathrm{s}$
plant command per step. Its primary action is normalized steering rate at the
fixed reverse-speed target. A separate two-action teleoperation mode adds speed
authority for manual debugging without changing the research action contract.
The reward combines progress with penalties for time, lateral and heading
error, articulation, and steering rate, plus explicit terminal terms. Every
component is returned in the diagnostic information and versioned in a strict
task configuration. Success, collision, jackknife, workspace exit, and timeout
are evaluated independently of reward.

The compact registered task remains the controller-comparison baseline. The
full-size interactive task is selected explicitly by separate vehicle, dock,
and task configurations. In that demonstration, the colored active loading
zone has no longitudinal side-wall collision polygons; complete-trailer
containment still uses the configured zone width. The rear boundary and
building façade outside the active door remain collidable. Its hard jackknife
event occurs at $60^\circ$, while the compact baseline terminates at
$45^\circ$.

\subsection{NMPC baseline}

\planned{} NMPC will predict the shared articulated-vehicle model and minimize
tracking, articulation, steering effort, and steering-rate variation while
enforcing dynamics, steering, steering-rate, and articulation constraints. The
primary comparison will retain fixed reverse speed and steering-rate authority.
Horizon, discretization, weights, warm start, solver tolerances, and numerical
failure handling will be recorded in versioned experiment configurations.

\subsection{PPO controller}

\planned{} PPO will use the same steering-rate authority and fixed reverse
speed. Training is expected to proceed from straight reverse stabilization to
curved tracking and randomized docking. Reward components, termination logic,
network architecture, training budget, and curriculum transitions will be
documented after implementation. A global-coordinate policy and removal of
preview information are candidate observation ablations, not established
experiments.

\subsection{Safety supervision}

\implemented{} Raw and supervised commands are represented separately so
interventions are observable. Steering-rate, steering-angle, speed, and
acceleration limits are enforced through the shared configuration. Chassis
polygons are checked against margin-adjusted dock walls, and hard collision and
articulation events are independent of learned reward.

\section{Experimental Protocol}

\planned{} Experiments will be identified by immutable IDs and linked to a
versioned configuration, Git state, dependency versions, vehicle and dock
fingerprints, seeds, output paths, validity status, and interpretation notes.
MuJoCo runs will additionally store the complete resolved tuning values because
those values are not part of the version-1 vehicle fingerprint.
Large checkpoints and trajectories will remain outside the tracked manuscript,
while compact metrics, summaries, and reproducible figures will be retained.

The primary comparison will evaluate PPO and NMPC on identical held-out starts,
target docks, geometry, plant limits, state availability, and safety definitions.
Methodological asymmetries will be stated explicitly: NMPC receives an explicit
prediction model, while PPO receives training experience and a selected
observation representation. Offline training/tuning cost will be separated
from online control computation.

Planned metrics are:
\begin{itemize}
  \item docking success, collision, and jackknife rates;
  \item final longitudinal, lateral, and trailer-heading error;
  \item maximum articulation angle;
  \item maneuver time and path length;
  \item steering-rate magnitude and smoothness;
  \item controller computation time and numerical failures; and
  \item safety-supervisor intervention count.
\end{itemize}

After nominal performance is established, candidate robustness studies include
held-out start poses, hitch-to-axle distances, bay dimensions, clearance, state
noise, and model mismatch. Geometry variation alone will be described as
sensitivity, robustness, or generalization; the term \emph{ablation} will be
reserved for removing or modifying a defined component.

Transfer evaluation will proceed only after 2D/MuJoCo parity checks establish
matching frames, signs, dimensions, command timing, semantic sites, and event
definitions. Zero-shot transfer, domain-randomized transfer, and MuJoCo
fine-tuning will be reported as separate conditions if they are implemented.

\section{Development Status and Remaining Work}

Table~\ref{tab:status} is the manuscript-level project map. A component is
marked implemented only when corresponding repository code and tests exist.

\begin{table*}[ht]
  \centering
  \caption{Current implementation status.}
  \label{tab:status}
  \begin{tabular}{p{0.34\textwidth}p{0.14\textwidth}p{0.43\textwidth}}
    \toprule
    Component & Status & Evidence or next gate \\
    \midrule
    Proposal and initial literature set & \implemented & Submitted proposal and four local study papers \\
    Versioned vehicle and dock contract & \implemented & Specification and strict YAML configuration \\
    Shared state, command, and geometry & \implemented & Unit tests and randomized transform checks \\
    Provenance and conformance tooling & \implemented & Fingerprints and executable conformance report \\
    Command supervision & \implemented & Auditable clipping and violation tests \\
    2D kinematic propagation & \implemented & RK4, equation, limit, rigid-geometry, determinism, and convergence tests \\
    Paths, previews, collision, and goal logic & \planned & Geometry and deterministic scenario tests \\
    Gymnasium environment & \planned & API checker and seeded rollout tests \\
    Simple closed-loop tracker & \planned & Straight and curved reference tracking \\
    NMPC & \planned & Solver, feasibility, and timing validation \\
    PPO & \planned & Training, evaluation, and curriculum evidence \\
    MuJoCo vehicle, dock, and adapter & \implemented & Generated MJCF; neutral-pose, sign, timing, contact, and kinematic-agreement tests \\
    2D/MuJoCo simulator parity & \implemented & Direct forward/reverse plant trajectories; task-event equality remains planned \\
    Comparative experiments & \planned & Frozen held-out scenarios and saved results \\
    \bottomrule
  \end{tabular}
\end{table*}

The immediate development gate is shared reference-path, footprint, collision,
jackknife, and goal logic. These modules must be verified independently before
they are composed into the Gymnasium environment and used for task-event parity.

\section{Results}

No controller-training or comparative experimental results exist yet.

The current evidence is limited to engineering validation of the shared core
and simulator plants. The repository test suite contains 225 passing tests,
including the 2D kinematic and MuJoCo checks, with branch-aware coverage of
$92.70\%$ across the measured packages. Lint and strict
static type checks pass. The executable conformance report verifies
the neutral rigid-hitch constraint, tractor and trailer footprint areas,
integer 2D/MuJoCo timing ratios, nominal speed bounds, nominal trailer fit in
the configured dock, and semantic-site coverage. Kinematic checks exercise
forward and reverse propagation, constant steering, articulation response,
acceleration and steering limits, rigid reconstruction, deterministic replay,
and timestep convergence. Shared command sequences also compare the final 2D
and MuJoCo planar states within explicit test tolerances.

These engineering checks establish internal consistency and low-speed model
behavior; they do not establish real-vehicle physical fidelity or controller
performance.
Numerical controller results, plots, confidence intervals, and failure
categories will be added only from saved experiment IDs.

\section{Discussion and Limitations}

The implemented core reduces a major source of integration error by assigning
explicit meanings to poses, signs, dimensions, semantic sites, timing, and
controller commands before simulator development. Strict configuration loading
and fingerprints also make later geometry changes traceable. This is an
engineering foundation rather than a scientific result.

Several important limitations remain. The compact vehicle is not a validated
full-size highway model. MuJoCo masses, contact parameters, actuator limits,
and tuned gains are engineering values rather than identified real-vehicle
parameters. The selected $0.30\,\mathrm{rad/s}$ steering-rate limit is not
inherited from the two core papers. The implemented tests now compare MuJoCo
directly with the project 2D backend for low-speed plant trajectories, but they
do not establish parity for future collision, safety, docking, observation, or
termination events.
MuJoCo-only tuning is versioned in code but not yet included in the vehicle
configuration fingerprint.

Potentially publishable observations may emerge from geometry generalization,
controller sensitivity, or 2D-to-3D transfer, but these remain open research
questions. Novelty will require repeatable evidence and an updated literature
assessment, not merely behavior that is new within this repository.

\section{Conclusion}

The project currently has a frozen, executable cross-simulator contract for a
compact tractor--semitrailer and loading dock. Shared state, command, geometry,
configuration, provenance, safety-supervision, and conformance-validation code
is implemented and tested. A configuration-generated MuJoCo vehicle, dock, and
shared-state adapter now satisfy the model-side contract and low-speed
engineering checks. The next milestone is the two-dimensional kinematic
state-transition model and its quantitative validation, followed by direct
cross-backend parity. Conclusions about PPO, NMPC, robustness, or transfer are
deferred until reproducible experiments exist.


\bibliographystyle{IEEEtran}
\bibliography{references}

\end{document}
